\documentclass{amsart}
% For dealing with references we use the comment environment
\usepackage{verbatim}
\newenvironment{reference}{\comment}{\endcomment}
%\newenvironment{reference}{}{}
\newenvironment{slogan}{\comment}{\endcomment}
\newenvironment{history}{\comment}{\endcomment}

% For commutative diagrams we use Xy-pic
\usepackage[all]{xy}

% We use 2cell for 2-commutative diagrams.
\xyoption{2cell}
\UseAllTwocells

% We use multicol for the list of chapters between chapters
\usepackage{multicol}

% This is generally recommended for better output
\usepackage{lmodern}
\usepackage[T1]{fontenc}

% For cross-file-references

% Package for hypertext links:
\usepackage{hyperref}

% For any local file, say "hello.tex" you want to link to please

% Theorem environments.
%
\theoremstyle{plain}
\newtheorem{theorem}[subsection]{Theorem}
\newtheorem{proposition}[subsection]{Proposition}
\newtheorem{lemma}[subsection]{Lemma}

\theoremstyle{definition}
\newtheorem{definition}[subsection]{Definition}
\newtheorem{example}[subsection]{Example}
\newtheorem{exercise}[subsection]{Exercise}
\newtheorem{situation}[subsection]{Situation}

\theoremstyle{remark}
\newtheorem{remark}[subsection]{Remark}
\newtheorem{remarks}[subsection]{Remarks}

\numberwithin{equation}{subsection}

% Macros
%
\def\lim{\mathop{\mathrm{lim}}\nolimits}
\def\colim{\mathop{\mathrm{colim}}\nolimits}
\def\Spec{\mathop{\mathrm{Spec}}}
\def\Hom{\mathop{\mathrm{Hom}}\nolimits}
\def\Ext{\mathop{\mathrm{Ext}}\nolimits}
\def\SheafHom{\mathop{\mathcal{H}\!\mathit{om}}\nolimits}
\def\SheafExt{\mathop{\mathcal{E}\!\mathit{xt}}\nolimits}
\def\Sch{\mathit{Sch}}
\def\Mor{\mathop{\mathrm{Mor}}\nolimits}
\def\Ob{\mathop{\mathrm{Ob}}\nolimits}
\def\Sh{\mathop{\mathit{Sh}}\nolimits}
\def\NL{\mathop{N\!L}\nolimits}
\def\CH{\mathop{\mathrm{CH}}\nolimits}
\def\proetale{{pro\text{-}\acute{e}tale}}
\def\etale{{\acute{e}tale}}
\def\QCoh{\mathit{QCoh}}
\def\Ker{\mathop{\mathrm{Ker}}}
\def\Im{\mathop{\mathrm{Im}}}
\def\Coker{\mathop{\mathrm{Coker}}}
\def\Coim{\mathop{\mathrm{Coim}}}

% Boxtimes
%
\DeclareMathSymbol{\boxtimes}{\mathbin}{AMSa}{"02}

%
% Macros for moduli stacks/spaces
%
\def\QCohstack{\mathcal{QC}\!\mathit{oh}}
\def\Cohstack{\mathcal{C}\!\mathit{oh}}
\def\Spacesstack{\mathcal{S}\!\mathit{paces}}
\def\Quotfunctor{\mathrm{Quot}}
\def\Hilbfunctor{\mathrm{Hilb}}
\def\Curvesstack{\mathcal{C}\!\mathit{urves}}
\def\Polarizedstack{\mathcal{P}\!\mathit{olarized}}
\def\Complexesstack{\mathcal{C}\!\mathit{omplexes}}
% \Pic is the operator that assigns to X its picard group, usage \Pic(X)
% \Picardstack_{X/B} denotes the Picard stack of X over B
% \Picardfunctor_{X/B} denotes the Picard functor of X over B
\def\Pic{\mathop{\mathrm{Pic}}\nolimits}
\def\Picardstack{\mathcal{P}\!\mathit{ic}}
\def\Picardfunctor{\mathrm{Pic}}
\def\Deformationcategory{\mathcal{D}\!\mathit{ef}}

\setlength{\emergencystretch}{3em}
\begin{document}
\title[Stacks Project, Tag 058X]{1273 --- Kaplansky decomposition of direct summands of sums of countably generated modules over arbitrary rings}
\maketitle
\noindent Copyright (C) 2005--2025 Johan de Jong. Stacks Project authors. Source: \href{https://stacks.math.columbia.edu/tag/058X}{Tag 058X}.

\noindent Editorial difficulty note (not author text). Difficulty H2: The author constructs a transfinite filtration compatible with two projections, using countable matrix closure along successive anti-diagonals to ensure every successor stage is a sum of original summands. This arbitrary-module structural decomposition requires substantive transfinite construction beyond qualification homological algebra..

\noindent Editorial provenance note (not author text). History: excerpt from revision \texttt{a0444\allowbreak 6e57e\allowbreak c1fbc\allowbreak 252a8\allowbreak 71afc\allowbreak ec775\allowbreak 2fb28\allowbreak 07b14}, algebra.tex, lines 20469--20566. Reference presentation was adapted; original-author context and core support blocks were retained or added. The mathematical wording is unchanged. Licensed under GNU FDL 1.2 or later; no invariant sections or cover texts. See ../licenses/COPYING. Context blocks reproduce author definitions and supporting results; they are not additional counted problems.

\begin{definition}
\label{definition-devissage}
Let $M$ be an $R$-module.  A {\it direct sum d\'evissage} of $M$ is a family
of submodules $(M_{\alpha})_{\alpha \in S}$, indexed by an ordinal $S$ and
increasing (with respect to inclusion), such that:
\begin{enumerate}
\item[(0)] $M_0 = 0$;
\item[(1)] $M = \bigcup_{\alpha} M_{\alpha}$;
\item[(2)] if $\alpha \in S$ is a limit ordinal, then $M_{\alpha} =
\bigcup_{\beta < \alpha} M_{\beta}$;
\item[(3)] if $\alpha + 1 \in S$, then $M_{\alpha}$ is a direct summand of
$M_{\alpha + 1}$.
\end{enumerate}
If moreover
\begin{enumerate}
\item[(4)] $M_{\alpha + 1}/M_{\alpha}$ is countably generated for
$\alpha + 1 \in S$,
\end{enumerate}
then $(M_{\alpha})_{\alpha \in S}$ is called a {\it Kaplansky d\'evissage}
of $M$.
\end{definition}

\begin{lemma}
\label{lemma-direct-sum-devissage}
Let $M$ be an $R$-module.  If $(M_{\alpha})_{\alpha \in S}$ is a direct sum
d\'evissage of $M$, then
$M \cong \bigoplus_{\alpha + 1 \in S} M_{\alpha + 1}/M_{\alpha}$.
\end{lemma}

\begin{proof}
By property (3) of a direct sum d\'evissage, there is an inclusion
$M_{\alpha + 1}/M_{\alpha} \to M$ for each $\alpha \in S$.  Consider the
map
$$
f : \bigoplus\nolimits_{\alpha + 1\in S} M_{\alpha + 1}/M_{\alpha} \to M
$$
given by the sum of these inclusions.
Further consider the restrictions
$$
f_{\beta} :
\bigoplus\nolimits_{\alpha + 1 \leq \beta} M_{\alpha + 1}/M_{\alpha}
\longrightarrow
M
$$
for $\beta\in S$. Transfinite induction on $S$ shows that the image of
$f_{\beta}$ is $M_{\beta}$. For $\beta=0$ this is true by $(0)$. If $\beta+1$
is a successor ordinal and it is true for $\beta$, then it is true for 
$\beta + 1$ by (3).  And if $\beta$ is a limit ordinal and it is true for
$\alpha < \beta$, then it is true for $\beta$ by (2). Hence $f$ is surjective
by (1).  

\medskip\noindent
Transfinite induction on $S$ also shows that the restrictions $f_{\beta}$
are injective. For $\beta = 0$ it is true. If $\beta+1$ is a
successor ordinal and $f_{\beta}$ is injective, then let $x$ be in the kernel
and write $x = (x_{\alpha + 1})_{\alpha + 1 \leq \beta + 1}$ in terms of its
components $x_{\alpha + 1} \in M_{\alpha + 1}/M_{\alpha}$.  By property (3) and
the fact that the image of $f_{\beta}$ is $M_{\beta}$ both
$(x_{\alpha + 1})_{\alpha + 1 \leq \beta}$ and $x_{\beta + 1}$ map to $0$.
Hence $x_{\beta+1} = 0$ and, by the assumption that the restriction
$f_{\beta}$ is injective also $x_{\alpha + 1} = 0$
for every $\alpha + 1 \leq \beta$.  So $x = 0$ and $f_{\beta+1}$ is injective.
If $\beta$ is a limit ordinal consider an element $x$ of the kernel. Then $x$
is already contained in the domain of $f_{\alpha}$ for some $\alpha < \beta$.
Thus $x = 0$ which finishes the induction.
We conclude that $f$ is injective since $f_{\beta}$ is for each $\beta \in S$.
\end{proof}


\begin{theorem}
\label{theorem-kaplansky-direct-sum}
Suppose $M$ is a direct sum of countably generated $R$-modules.  If $P$ is a
direct summand of $M$, then $P$ is also a direct sum of countably generated
$R$-modules.
\end{theorem}

\begin{proof}
Write $M = P \oplus Q$.  We are going to construct a Kaplansky d\'evissage
$(M_{\alpha})_{\alpha \in S}$ of $M$ which, in addition to the defining
properties (0)-(4), satisfies:
\begin{enumerate}
\item[(5)] Each $M_{\alpha}$ is a direct summand of $M$;
\item[(6)] $M_{\alpha} = P_{\alpha} \oplus Q_{\alpha}$, where $P_{\alpha} =P
\cap M_{\alpha}$ and $Q_\alpha = Q \cap M_{\alpha}$.
\end{enumerate}
(Note: if properties (0)-(2) hold, then in fact property (3) is equivalent to
property (5).)

\medskip\noindent
To see how this implies the theorem, it is enough to show that
$(P_{\alpha})_{\alpha \in S}$ forms a Kaplansky d\'evissage of $P$.  Properties
(0), (1), and (2) are clear.  By (5) and (6) for $(M_{\alpha})$, each
$P_{\alpha}$ is a direct summand of $M$.  Since $P_{\alpha} \subset P_{\alpha +
1}$, this implies $P_{\alpha}$ is a direct summand of $P_{\alpha + 1}$; hence
(3) holds for $(P_{\alpha})$.  For (4), note that
$$
M_{\alpha + 1}/M_{\alpha} \cong P_{\alpha + 1}/P_{\alpha} \oplus
Q_{\alpha + 1}/Q_{\alpha},
$$
so $P_{\alpha + 1}/P_{\alpha}$ is countably generated because this is true of
$M_{\alpha + 1}/M_{\alpha}$.

\medskip\noindent
It remains to construct the $M_{\alpha}$.  Write $M = \bigoplus_{i \in I} N_i$
where each $N_i$ is a countably generated $R$-module.  Choose a well-ordering
of $I$.  By transfinite recursion we are going to define an increasing family
of submodules $M_{\alpha}$ of $M$, one for each ordinal $\alpha$, such that
$M_{\alpha}$ is a direct sum of some subset of the $N_i$.

\medskip\noindent
For $\alpha = 0$ let $M_{0} = 0$.  If $\alpha$ is a limit ordinal and
$M_{\beta}$ has been defined for all $\beta < \alpha$, then define $M_{\alpha}
= \bigcup_{\beta < \alpha} M_{\beta}$.  Since each $M_{\beta}$ for $\beta <
\alpha$ is a direct sum of a subset of the $N_i$, the same will be true of
$M_{\alpha}$.  If $\alpha + 1$ is a successor ordinal and $M_{\alpha}$ has been
defined, then define $M_{\alpha + 1}$ as follows.  If $M_{\alpha} = M$, then let
$M_{\alpha + 1} = M$.  If not, choose the smallest $j \in I$ such that $N_j$ is
not contained in $M_{\alpha}$.  We will construct an infinite matrix $(x_{mn}),
m, n = 1, 2, 3, \ldots$ such that:
\begin{enumerate}
\item $N_j$ is contained in the submodule of $M$ generated by the entries
$x_{mn}$;
\item if we write any entry $x_{k\ell}$ in terms of its $P$- and
$Q$-components, $x_{k\ell} = y_{k\ell} + z_{k\ell}$, then the matrix $(x_{mn})$
contains a set of generators for each $N_i$ for which $y_{k\ell}$ or
$z_{k\ell}$ has nonzero component.
\end{enumerate}
Then we define $M_{\alpha + 1}$ to be the submodule of $M$ generated by
$M_{\alpha}$ and all $x_{mn}$; by property (2) of the matrix $(x_{mn})$,
$M_{\alpha + 1}$ will be a direct sum of some subset of the $N_i$.
To construct the matrix $(x_{mn})$, let $x_{11}, x_{12}, x_{13}, \ldots$
be a countable set of generators for $N_j$. Then if
$x_{11} = y_{11} + z_{11}$ is the decomposition into $P$- and
$Q$-components, let $x_{21}, x_{22}, x_{23}, \ldots$ be a countable
set of generators for the sum of the $N_i$ for which $y_{11}$ or $z_{11}$ have
nonzero component.  Repeat this process on $x_{12}$ to get elements $x_{31},
x_{32}, \ldots$, the third row of our matrix.  Repeat on $x_{21}$ to get the
fourth row, on $x_{13}$ to get the fifth, and so on, going down along
successive anti-diagonals as indicated below:
$$
\left(
\vcenter{
\xymatrix@R=2mm@C=2mm{
x_{11} & x_{12} \ar[dl] & x_{13} \ar[dl] & x_{14} \ar[dl] & \ldots  \\
x_{21} & x_{22} \ar[dl] & x_{23} \ar[dl] & \ldots  \\
x_{31} & x_{32} \ar[dl] & \ldots \\
x_{41} & \ldots \\
\ldots
}
}
\right).
$$

\medskip\noindent
Transfinite induction on $I$ (using the fact that we constructed
$M_{\alpha + 1}$
to contain $N_j$ for the smallest $j$ such that $N_j$ is not contained in
$M_{\alpha}$) shows that for each $i \in I$, $N_i$ is contained in some
$M_{\alpha}$.  Thus, there is some large enough ordinal $S$ satisfying: for
each $i \in I$ there is $\alpha \in S$ such that $N_i$ is contained in
$M_{\alpha}$.  This means $(M_{\alpha})_{\alpha \in S}$ satisfies property (1)
of a Kaplansky d\'evissage of $M$.  The family $(M_{\alpha})_{\alpha \in S}$
moreover satisfies the other defining properties, and also (5) and (6) above:
properties (0), (2), (4), and (6) are clear by construction;  property (5) is
true because each $M_{\alpha}$ is by construction a direct sum of some $N_i$;
and (3) is implied by (5) and the fact that $M_{\alpha} \subset M_{\alpha + 1}$.
\end{proof}
\end{document}
